\section{Bounded-density compression with factored Gaussian reconstruction}
\label{app:factor-code}

The all-scale code in Appendix~\ref{app:all-scales-quantization} can be
sharpened in two further ways. First, the actual conditioned posterior has
a bounded density relative to the raw Gaussian. Propagating this stronger
property through the virtual computation removes the factor $1/\eta$
from the active-direction cutoff. This argument is specific to bounded
density; it does not upgrade the compression theorem for arbitrary laws
having only a relative-entropy bound. Second, the active part of each
reconstructed matrix already has small rank, and its fresh Gaussian part
can be replaced by a small-rank Gaussian factorization. Coding the factors
then gives a sufficient sample order $\widetilde O(C^5L^{12})$ at every
$0<t\le1$. The support remains exactly $\mathcal R_{N,c}$. This is a supporting
code: the entropy-metric and Gaussian-weighted constructions in the
following appendices improve its fifth-power width dependence to the
headline quartic bound.

We give the density and factorization arguments below. We use the
whole-layer active-direction and whitened-sampling estimates stated in
Appendix~\ref{app:all-scales-quantization}, with their original proofs
in~\cite{all-scales}. As in that appendix, every comparison coupling is
allowed to depend on the input, whereas the candidate law is not. The
argument and the improved exponents are conventional proofs and are not
Lean formalized.

\subsection{Explicit budgets and learning statement}

Use the product-cylinder cost $D_*=D_*^{\rm prod}$ from
\eqref{eq:product-cylinder-cost}. For $0<\varepsilon,\delta<1$, define
\begin{align}
 \beta&=\varepsilon/4,&
 B_v&=\max\{8,8[D_*+\log(1/\beta)]\},&
 K&=D_*+1,\notag\\
 r&=\min\{1/6,c/\sqrt{1+(L-1)B_v}\},&
 \alpha&=\varepsilon/12,&
 e&=\frac{\alpha r^2}{1000L},\qquad
 \eta=\frac{\alpha}{100L}.
 \label{eq:factor-posterior-budgets}
\end{align}
Set
\begin{align}
 q&=\left\lceil\frac{2K+\log(1/\eta)}e\right\rceil,&
 D_0&=[(d+1)q+2]\log2,&
 V&=8K+2q,\notag\\
 B_H^2&=4V/e,&
 K_0&=(4K+2\log2)/\eta,&
 B_1^2&=4[4K+2(d+1)\log2]/e.
 \label{eq:factor-active-budgets}
\end{align}
Let $m$ be the least integer no smaller than each of
\begin{align}
 &10000(V+q+1)/e,\notag\\
 &64(q+1)\log(4(q+1)/\eta),\label{eq:factor-source-width}\\
 &512K_0r^{-2}[d\log2+\log(2/\eta)+1].\notag
\end{align}
The factor rank and the small-width cutoff are
\begin{equation}\label{eq:factor-rank-cutoff}
 p=\lceil1000/e\rceil,\qquad k=q+p,\qquad
 n_0=\max\{1,\lceil LD_0/K\rceil\}.
\end{equation}
In particular, the cutoff $n_0$ need not be as large as $m$.
Define factor bounds and parameter counts by
\begin{align}
 B_g^2&=4p/e,\qquad A_f=\max\{B_H,B_g,1\},\notag\\
 B_f&=\max\{2\sqrt{q+1}/m,\sqrt{2/\eta},1\},\notag\\
 P_{\rm small}&=n_0(d+1)+(L-2)n_0^2+n_0,\notag\\
 P_{\rm fac}&=m(d+1)+2(L-2)mk+m,\notag\\
 B_{\rm small}^2&=\frac{4(4n_0K+2P_{\rm small}\log2)}\alpha,\notag\\
 B_{\rm code}&=1+\max\{B_{\rm small},B_1,A_f,B_f,
                                      \sqrt{K_0}/m,1\}.
 \label{eq:factor-code-bounds}
\end{align}
Finally, choose the common scalar mesh
\begin{equation}\label{eq:factor-mesh}
 \zeta=\min\left\{1,
 \frac{\sqrt e}{4(d+1)},\frac{\sqrt e}{4m},
 \frac{\sqrt e}{4mk(A_f+B_f+1)},
 \frac{\sqrt e}{4(n_0+d+1)}\right\},
\end{equation}
and put
\begin{equation}\label{eq:factor-code-budget}
 J_{\rm fac}=\log2+
 \max\{P_{\rm small},P_{\rm fac}\}
               \log(3+4B_{\rm code}/\zeta).
\end{equation}
All these quantities are independent of learner width and Gaussian scale.

\begin{theorem}[All-scale learning by a factored code]
\label{thm:factor-all-scales}
Under the Boolean-teacher convention of the main theorem, the condition
\begin{equation}\label{eq:factor-sample-budget}
 n\ge\frac{16}{\varepsilon}
             [J_{\rm fac}+\log(1/\delta)]
\end{equation}
implies, with probability at least $1-\delta$ over iid training inputs,
\[
 \E_{w\sim Q_{S,N,t,c},\,X\sim D}
       \mathbf1\{f(X)z_w(X)\le1/4\}\le\varepsilon
\]
simultaneously for every $N\ge C$ and every $0<t\le1$.
The same conclusion holds if the posterior conditioning additionally
requires the label-independent global margin $|z_w(x)|>1/2$ for every
input $x$.
No population-merging assumption or additional internal-noise condition
is imposed.
\end{theorem}

\begin{corollary}[Improved width and depth powers]
\label{cor:factor-powers}
For fixed $d,c,\varepsilon,\delta$, let $F_0=1+\log^2(2+C)$.
A sufficient sample count in Theorem~\ref{thm:factor-all-scales} has order
\begin{equation}\label{eq:factor-powers}
 O\!\left(C^5L^{12}F_0^5\log(2+CL)\right)
       =\widetilde O(C^5L^{12}).
\end{equation}
At fixed teacher width this is $O(L^{12}\log(2+L))$; at fixed depth it is
$O(C^5\log^{11}(2+C))$.
\end{corollary}

\subsection{A bounded-density invariant for unprocessed weights}

The following lemma strengthens the entropy bookkeeping at the point
where an active rank must hold with high probability. A measure inequality
means domination on every measurable set.

\begin{lemma}[Conditional density and finite-message propagation]
\label{lem:factor-density}
Suppose that
\[
 Q_{A,X,Y}\le e^b(\nu_A\otimes P_X\otimes P_Y),
\]
where $\nu_A,P_X,P_Y$ are probability reference laws. Write
$u=dQ_{A,Y}/d(\nu_A\otimes P_Y)$. Then
\[
 Q\{u(A,Y)<\eta\}\le\eta,
\]
and, outside this event,
\begin{equation}\label{eq:factor-conditional-entropy}
 D_{\rm KL}(Q_{X\mid A,Y}\Vert P_X)
       \le b+\log(1/\eta).
\end{equation}
Suppose also that $J=j(A,Y)$ takes values in a fixed finite alphabet
$\mathcal J$, that $H_0=h(A,J,X)$, and that a new prefix $A'$ is generated
from $(A,J,H_0)$ through a kernel whose fresh randomness has no further
dependence on $Y$. There is a probability reference law $\nu_{A'}$ such
that
\begin{equation}\label{eq:factor-density-propagation}
 Q_{A',Y}\le e^{b+\log|\mathcal J|}
                            (\nu_{A'}\otimes P_Y).
\end{equation}
The original domination also implies
\[
 \E_{Q_A}D_{\rm KL}(Q_{X,Y\mid A}\Vert P_X\otimes P_Y)\le b.
\]
\end{lemma}
\begin{proof}
Let $F\le e^b$ be the joint density relative to
$\nu_A\otimes P_X\otimes P_Y$. Disintegration gives conditional density
$F(a,x,y)/u(a,y)$ for $Q_{X\mid a,y}$ wherever $u>0$.
Thus its relative entropy is at most $b-\log u(a,y)$.
Also
\[
 Q\{u<\eta\}=
 \int_{\{u<\eta\}}u\,d(\nu_A\otimes P_Y)\le\eta.
\]
This is an unconditional probability bound under the canonical joint law;
it need not hold separately at every fixed prefix.

For propagation, augment the reference law by an independent uniformly
distributed message in $\mathcal J$. The actual augmented law satisfies
\[
 Q_{A,J,X,Y}\le e^b|\mathcal J|
  (\nu_A\otimes\operatorname{Unif}(\mathcal J)
                       \otimes P_X\otimes P_Y).
\]
Apply the retained-data map and the same prefix kernel to both sides.
Under the reference law, its inputs and fresh randomness are independent
of $Y$, so its pushforward factors as $\nu_{A'}\otimes P_Y$.
This proves \eqref{eq:factor-density-propagation}. The last assertion
follows from the KL chain rule and the nonnegativity of
$D_{\rm KL}(Q_A\Vert\nu_A)$.
\end{proof}



Suppose initially that a law $Q$ on original weights satisfies
\begin{equation}\label{eq:factor-initial-density}
 Q\le\exp(NK/t^2)\,\mu_{N,t}.
\end{equation}
The virtual prefix records the constructed network, not just its value at
one selected input. Clipping and rounding the first layer are deterministic
channels applied to that original block and initialize the domination
invariant with exponent $NK/t^2$. At an internal stage, let $X$ be the
current original block and $Y$ the remaining original blocks. Induction
with Lemma~\ref{lem:factor-density} gives exponent
\[
 b_s\le NK/t^2+sD_0\le2NK/t^2\qquad(N>n_0).
\]
At that stage abort if its marginal density $u(A,Y)$ is below $\eta$.
Otherwise \eqref{eq:factor-conditional-entropy} and the imported
whole-layer active-direction lemma, applied with tolerance $e$, give
\[
 \dim U\le\frac{t^2}{Ne}
             [b_s+\log(1/\eta)]
       \le\frac{2K+\log(1/\eta)}e\le q.
\]
The active query list is selected from the conditional law given $(A,Y)$,
so it is a function of $(A,Y)$. Its basis is a function of the prefix and
the list, using fixed Gram--Schmidt conventions. The selected list and an
abort symbol admit a fixed alphabet of log-cardinality at most $D_0$.
For reference messages never chosen by the actual selection rule, extend
the basis and transition kernel by fixed measurable default rules. This
makes the uniform-message reference construction in the lemma defined
on its entire domain.

All abort paths remain in the canonical probability law. On an abort,
use a fixed default prefix, retain an unsuccessful flag, and use default
transitions thereafter. Do not condition any intermediate law on success.
The density-test abort is included in the finite message; later aborts
arising from fresh randomness are part of the prefix channel. Consequently
the invariant applies on the full law, and the unsuccessful paths are
discarded only in the final approximation probability estimate.

The domination invariant also provides the old averaged entropy bound.
Thus the imported Gaussian entropy-moment estimate still gives
\begin{equation}\label{eq:factor-active-moment}
 \E\left[\norm{H_0}_F^2/N;\ \text{active success}\right]
       \le4t^2b_s/N+2t^2q\log2\le V.
\end{equation}
In particular, the definitions of $V,B_H,B_1$, and $K_0$ need no extra
factor of $1/\eta$. The improvement in $q$ uses the density hypothesis
\eqref{eq:factor-initial-density}; an arbitrary KL bound alone would not
justify it.

For comparison, using this new rank cutoff with the dense sequential code
of Appendix~\ref{app:all-scales-quantization} already gives
$O(C^6L^{15}F_0^6\log(2+CL))$ samples. The factor construction below saves
additional powers of both resources.

\subsection{A small-rank replacement for the Gaussian field}

We recall the source-sampling notation needed for the new step.
The source activations obey $|h_j(x)|\le1$. If $U$ has $\ell\le q$
orthonormal columns, put $u_j=\sqrt N U_{j,:}$. The half-uniform,
half-leverage law gives weights $a_j=1/(Np_j)\le2$ and
$a_j\norm{u_j}^2\le2\ell$. Draw $m$ source indices $J_b$, and define
\[
 \widehat G=\frac1m\sum_ba_{J_b}u_{J_b}u_{J_b}^{\top},\qquad
 D_{b,:}=\sqrt{a_{J_b}/m}\,u_{J_b}^{\top}\widehat G^{-1/2}.
\]
On the good Gram event, $\widehat G\succeq I/2$ and $D^\top D=I$.
For $\ell=0$, the active matrices are empty and the following projection
is the identity. Write
\[
 P=I-DD^\top,\qquad
 c_b(x)=\sqrt{a_{J_b}/m}\,h_{J_b}(x),\qquad
 \widehat C(x,x)=\norm{Pc(x)}^2\le2.
\]
The ideal field in the imported whitened-sampling estimate has active mean
$H_0\psi(x)$ and independent Gaussian rows of variance $t^2C(x,x)$.
The original reconstruction has mean $H_0\widehat v(x)$ and variance
$t^2\widehat C(x,x)$.

\begin{lemma}[Factored Gaussian reconstruction]
\label{lem:factor-gaussian}
Conditional on the source samples and retained active data, draw an
independent standard Gaussian matrix $R\in\R^{p\times m}$ and independent
standard Gaussian row vectors $g_i\in\R^p$, for all virtual target rows.
Replace the reconstructed coefficients by
\begin{equation}\label{eq:factor-coefficients}
 \widetilde W_{ib}=
   \frac{a_{J_b}}m H_{0,i}\mathbin\cdot u_{J_b}
   +t\sqrt{\frac{a_{J_b}}{mp}}\,(g_i^\top RP)_b.
\end{equation}
This coefficient law does not depend on the input. At every fixed input
it can be coupled to the ideal field so that the expected squared RMS
evaluation error, with the good Gram indicator, is at most
\begin{equation}\label{eq:factor-reconstruction-error}
 \frac{4\norm{H_0}_F^2/N+32(q+1)}m+\frac8p.
\end{equation}
The coefficient matrix in \eqref{eq:factor-coefficients} has rank at most
$\ell+p$.
\end{lemma}
\begin{proof}
Fix the input and abbreviate $\sigma^2=\widehat C(x,x)$.
Conditional on $R$, the new Gaussian evaluations are independent across
rows with variance $t^2s^2$, where
\[
 s=\norm{RPc(x)}/\sqrt p.
\]
If $\sigma>0$, then $s^2/\sigma^2$ has the law of a chi-square variable
with $p$ degrees of freedom divided by $p$. Rationalizing the square root
therefore gives
\begin{equation}\label{eq:factor-sd-error}
 \E_R(s-\sigma)^2
 \le\sigma^2\E(\chi_p^2/p-1)^2
 =2\sigma^2/p\le4/p.
\end{equation}
For $\sigma=0$ the statement is immediate.

To make the coupling explicit, draw $R$ and independent standard normals
$Z_i$, independent of $R$. Use Gaussian row evaluations $t\sigma Z_i$
for the old reconstruction and $tsZ_i$ for the new one. Disintegrating
$g_i$ conditional on its projection along $RPc(x)$ realizes the latter
evaluation while keeping its full conditional marginal standard Gaussian.
Thus the coefficient law is exactly the law in
\eqref{eq:factor-coefficients}. Disintegration likewise extends the old
Gaussian evaluation to its coefficient field. The degenerate zero-variance
case is handled by independent completion. Only this comparison coupling,
not either field marginal, depends on $x$.

The old whitened reconstruction estimate is
$[2\norm{H_0}_F^2/N+16(q+1)]/m$. Couple through that reconstruction and
apply the squared triangle inequality with
\eqref{eq:factor-sd-error}. This yields
\eqref{eq:factor-reconstruction-error}. Finally the active term factors
through $\R^\ell$ and the Gaussian term through $\R^p$, proving the
rank claim.
\end{proof}

Integrating \eqref{eq:factor-reconstruction-error} with the success
indicator and using \eqref{eq:factor-active-moment}, the first line of
\eqref{eq:factor-source-width}, and $p\ge1000/e$, bounds this error by
\[
 \frac{4V+32(q+1)}m+\frac8p<e/8.
\]
The comparison requires no simultaneous coupling of all inputs.
Conditional independence of rows given $R$ suffices; independence of the
new rows after averaging over $R$ is not asserted or needed.

\subsection{Bounded factors and their sequential rounding}

At each internal stage first clip active rows of norm greater than $B_H$
to zero, exactly as in the imported construction. Equation
\eqref{eq:factor-active-moment} and bounded activations give squared RMS
activation cost at most $4V/B_H^2=e$. Apply
Lemma~\ref{lem:factor-gaussian} to this clipped matrix.

Abort if $\norm R_F^2>pm/\eta$. Since $\E\norm R_F^2=pm$, this event
has conditional probability at most $\eta$. Also replace each $g_i$ by
zero when $\norm{g_i}^2>B_g^2=4p/e$. This row replacement has probability
at most $e/4$ and hence expected squared activation change at most $e$:
two tanh values differ by at most two. All virtual rows are clipped before
any later row selection, including selections depending on future weights.

On success, write the resulting coefficients as $A_iB$, with zero padding
to factor dimension $k=q+p$, where
\[
 A_i=(H_{0,i},g_i),\qquad
 B_{{\rm active},ab}=a_{J_b}u_{J_b,a}/m,
\]
\[
 B_{{\rm Gaussian},ab}
       =t\sqrt{a_{J_b}/(mp)}\,(RP)_{ab}.
\]
Every left factor entry has absolute value at most $A_f$. Every right
factor entry has absolute value at most $B_f$, since
\[
 |a_{J_b}u_{J_b,a}/m|\le2\sqrt{q+1}/m,
\qquad
 t\sqrt{\frac{a_{J_b}}{mp}}|(RP)_{ab}|
 \le\sqrt{\frac2{mp}}\norm R_F\le\sqrt{2/\eta}.
\]

Round both factor arrays to the common scalar mesh \eqref{eq:factor-mesh}
before computing the next virtual prefix. For each bounded incoming state,
\begin{equation}\label{eq:factor-rounding-injection}
 |(A_{i,\rm round}B_{\rm round}-A_iB)h|
 \le mk(A_f+B_f+1)\zeta\le\sqrt e/4.
\end{equation}
Here each product changes by at most
$(A_f+B_f+1)\zeta$, because $\zeta\le1$.
Round first-layer coefficients and reconstructed readout coefficients
directly. Their preactivation errors are at most $(d+1)\zeta$ and
$m\zeta$, respectively, again at most $\sqrt e/4$.
The rounded candidate is an ordinary depth-$L$ tanh network: its internal
matrix is the product of two rounded factor arrays. Factorization adds
no activation layer.

All source samples, $R$, and Gaussian row factors are fresh randomness
conditional on the prefix, query list, and retained active matrix.
Clipping, factor multiplication, and rounding are deterministic subsequent
operations. They therefore form a channel of the precise kind required
by Lemma~\ref{lem:factor-density}. The reference-prefix law need not be
Gaussian or have a density on its own coordinate space. The comparison
coupling does not replace this canonical channel in the density induction.

\subsection{The disturbance budget and the finite class}

Consider first $N>n_0$ and a law satisfying
\eqref{eq:factor-initial-density}. The first layer is clipped at the
physical standardized row bound $B_1$ and then rounded. Its expected
squared RMS activation injection is less than $3e$ by the imported
first-layer estimate and the squared triangle inequality.
At an internal stage the five contributions are active-direction
transport, active-row clipping, reconstruction, Gaussian-row clipping,
and factor rounding. Their expected squared RMS bounds, with success
indicators, are respectively
\[
 2e,\quad e,\quad e/8,\quad e,\quad e/16.
\]
Another squared triangle inequality bounds their sum by
$5(2+1+1/8+1+1/16)e<21e$. In particular the conservative bound $41e$
from Appendix~\ref{app:all-scales-quantization} remains available.

The good Gram probability follows from the second line of
\eqref{eq:factor-source-width}. Use the original readout construction:
truncate at $N\norm v^2\le K_0$, sample source indices with probabilities
proportional to $|v_i|$, and use signed coefficients of magnitude
$\norm v_1/m\le\sqrt{K_0}/m$. The last line of
\eqref{eq:factor-source-width} gives output sampling error at most $r/16$
uniformly over the finite input domain except with probability $\eta$.
The energy truncation also fails with probability at most $\eta$.
Including density-test, Gram, and shared-sketch aborts at internal stages,
the total abort probability is at most $4L\eta$.

For each fixed input, glue the conditional transports as in
Appendix~\ref{app:all-scales-quantization}: retain the full original-weight
marginal $Q$ and the canonical outgoing prefix/future-weight marginal at
each transition. Extend the transition by the previous conditional law of
past original weights. This preserves every earlier computed prefix.
The final candidate output is exactly the original network under the
postactivation and readout injections defined by its rounded virtual
computation. No robustness of an intermediate candidate is assumed.

On the non-aborted paths, Markov's inequality gives
\[
 \Prb\left\{\sum_{s<L}\norm{\eta^s}^2/N>r^2/16,
                         \ \text{no abort}\right\}
 \le\frac{656Le}{r^2}=0.656\alpha.
\]
Readout rounding adds at most $\sqrt e/4\le r/64$. Thus, outside probability
at most $0.656\alpha+4L\eta<\alpha$, the total disturbance satisfies
\[
 \norm\eta_*\le
 \sqrt{r^2/16+(r/16+r/64)^2}
       =\frac{\sqrt{281}}{64}r<r/2.
\]

The construction above does not require $m\le N$. Source sampling is with
replacement, and duplicated indices give duplicated stored neurons. The
Gram lemma also allows this: on its good event $D$ has orthonormal columns,
and $\ell\le q<m$. Hence the only width cutoff needed for this branch is
$N>n_0$, which pays for the accumulated message exponents.

For $N\le n_0$, the original density bound implies a KL bound
$NK/t^2$. The Gaussian entropy-moment inequality gives, for all physical
standardized coordinates together,
\[
 \E_Q\norm{tZ_{\rm all}}^2
       \le4n_0K+2P_{\rm small}\log2.
\]
With probability at least $1-\alpha/4$, every actual coefficient is at
most $B_{\rm small}$ in absolute value. On that event round the original
network directly; otherwise use a zero-output default. The bounded rounded
activations give the exact injection bound
\[
 \norm\eta_*^2
 \le\zeta^2[(d+1)^2+(L-1)N^2]
 \le\zeta^2L(n_0+d+1)^2\le Le/16<r^2/4.
\]
Pad these small networks only to width $n_0$.

The large branch keeps the $m$ selected rows at each layer. It stores
$m(d+1)$ first-layer scalars, $2mk$ factor scalars for each internal
matrix, and $m$ readout scalars. The small branch stores at most
$P_{\rm small}$ scalars. Every rounded scalar in either branch belongs
to a fixed alphabet of at most $3+4B_{\rm code}/\zeta$ values.
The union of these two fixed finite classes, denoted $\mathcal H_{\rm fac}$,
therefore satisfies
\begin{equation}\label{eq:factor-finite-class}
 \log|\mathcal H_{\rm fac}|\le J_{\rm fac}.
\end{equation}
In particular one does not count the small branch as a dense network of
width $m$.

We have proved that every law satisfying
\eqref{eq:factor-initial-density} admits a single candidate law
$\Gamma_Q$ on $\mathcal H_{\rm fac}$ with the following property: at each
fixed input it has a coupling with $w\sim Q$ and a disturbance such that
\begin{equation}\label{eq:factor-final-coupling}
 h(x)=z_w(x;\eta),\qquad \norm\eta_*<r/2
\end{equation}
except with probability at most $\alpha$. Both marginals are independent
of the input used to form the comparison coupling.

\subsection{Posterior comparison and polynomial accounting}

\begin{proof}[Proof of Theorem~\ref{thm:factor-all-scales}]
The product target cylinder gives raw conditioning probability at least
$\exp(-ND_*/t^2)$. The chi-square exponential moment for the readout,
combined with this denominator, gives
\[
 Q_{S,N,t,c}\{E_v>B_v\}
 \le\exp\left(-\frac{NB_v}{4t^2}
                  +\frac N2\log2+\frac{ND_*}{t^2}\right)
 \le\beta.
\]
Let $Q'$ be the posterior conditioned on $E_v\le B_v$. Unlike a generic
law controlled only in relative entropy, it satisfies the literal
density bound
\[
 \frac{dQ'}{d\mu_{N,t}}
 \le\frac{\exp(ND_*/t^2)}{1-\beta}
 \le\exp(NK/t^2).
\]
Every original network in its support is robust on the common radius $r$.
Consequently \eqref{eq:factor-final-coupling} implies clipped-score error
at most $1/16$ on its good event.

Use the ramp $\psi$ from the earlier statistical proof: it is one on
$(-\infty,1/3]$, zero on $[5/12,\infty)$, and affine in between.
Every training point has candidate ramp loss zero on the good coupling;
every original test-margin failure has candidate loss one there.
A relative Chernoff bound and a union bound over
$\mathcal H_{\rm fac}$ give, on one sample event of probability at least
$1-\delta$,
\[
 \E_D\psi(f(X)T(h(X)))
 \le\frac2n\sum_{i=1}^n\psi(f(X_i)T(h(X_i)))
       +\frac{8[J_{\rm fac}+\log(1/\delta)]}{n}
 \quad(h\in\mathcal H_{\rm fac}).
\]
Averaging this inequality under $\Gamma_{Q'}$, applying the two marginal
comparisons, and restoring the discarded readout tail bounds the desired
risk by
\[
 \beta+3\alpha+
       \frac{8[J_{\rm fac}+\log(1/\delta)]}{n}\le\varepsilon.
\]
The finite class and sample event do not depend on $N,t$, which gives the
simultaneous conclusion. Additional global clean-margin conditioning retains
the same target cylinder, density bound, and readout-tail argument, so the
same proof applies to that stated optional variant.
\end{proof}

\begin{proof}[Proof of Corollary~\ref{cor:factor-powers}]
The product-cylinder constants give, at fixed ancillary resources,
\[
 K,B_v=O(CLF_0),\qquad r^{-2}=O(CL^2F_0),\qquad
 e^{-1}=O(CL^3F_0),\qquad \eta^{-1}=O(L).
\]
The term $\log(1/\eta)$ in the rank bound is dominated by the target cost
$K$. Therefore
\[
 q,V=O(C^2L^4F_0^2),\qquad p=O(CL^3F_0),\qquad
 m=O(C^3L^7F_0^3),\qquad n_0=O(CL^4F_0).
\]
The Gram and readout terms in the maximum defining $m$, including their
logarithms, are dominated by its first term at this fixed-resource scale.
It follows that
\[
 P_{\rm fac}=O(C^5L^{12}F_0^5),\qquad
 P_{\rm small}=O(C^2L^9F_0^2).
\]
All coefficient bounds and inverse meshes are polynomial in $C,L,F_0$,
so $\log(3+4B_{\rm code}/\zeta)=O(\log(2+CL))$.
Substitution into \eqref{eq:factor-code-budget} and
\eqref{eq:factor-sample-budget} proves the claimed order and both special
cases.
\end{proof}

\begin{remark}[A conservative joint polynomial for the factored code]
\label{rem:factor-joint-polynomial}
Put
\[
 S=2+d+L+C+c+c^{-1}+\varepsilon^{-1}+\log(1/\delta).
\]
The target-cylinder estimates and
\eqref{eq:factor-posterior-budgets} give
\[
 K,B_v=O(S^3),\qquad r^{-2}=O(S^6),\qquad
 e^{-1}=O(S^8),\qquad \eta^{-1}=O(S^2).
\]
The new budgets consequently satisfy
\[
 q,V=O(S^{11}),\quad D_0=O(S^{12}),\quad
 m=O(S^{19}),\quad n_0=O(S^{13}),\quad p=O(S^8),
\]
\[
 P_{\rm fac}=O(S^{31}),\quad P_{\rm small}=O(S^{27}),\quad
 B_{\rm small}^2=O(S^{28}),\quad B_{\rm code}=O(S^{14}).
\]
For the mesh, $A_f=O(S^{10})$ and $B_f=O(S)$; the latter uses the
first lower bound on $m$ to bound $\sqrt{q+1}/m$ by a constant.
Thus $\zeta^{-1}=O(S^{44})$, where the factor-rounding denominator
dominates the other displayed choices. It follows that
\[
 J_{\rm fac}=O(S^{31}\log S),\qquad
 n=O(S^{32}\log S)\subseteq O(S^{33}).
\]
This degree-33 envelope is deliberately conservative. Like the refined
width and depth powers above, it is a consequence of the present
conventional argument and carries no Lean-certification claim.
\end{remark}
